\section{The Celebrity Problem}
\label{sec:sq-celebrity-problem}

\problemstatement{There are $n$ people at a party. A \textbf{celebrity}
is a person whom \textbf{everyone else knows}, but who \textbf{knows no
one else}. We have access to a function $\textit{knows}(a,b)$ that
returns whether person $a$ knows person $b$. Find the celebrity if one
exists (there is at most one), using as few calls to $\textit{knows}$ as
possible -- ideally $O(n)$, not the $O(n^2)$ a full pairwise check would
need.}

\begin{patternbox}
\emph{``Repeatedly eliminate one of two candidates based on a pairwise
comparison, narrowing down to a single survivor''} is a stack-driven
elimination tournament: push every person as a candidate, then repeatedly
pop two candidates and use a single $\textit{knows}$ call to eliminate
(rather than confirm) one of them, since the celebrity's defining
properties make it possible to always rule out at least one of any two
candidates with a single check.
\end{patternbox}

\subsection*{Understanding the problem carefully}

Take any two candidates $a$ and $b$, and ask $\textit{knows}(a,b)$.

\begin{itemize}
  \item If $a$ knows $b$: then $a$ cannot be the celebrity (a celebrity
        knows \emph{no one}), so $a$ is eliminated.
  \item If $a$ does not know $b$: then $b$ cannot be the celebrity (the
        celebrity must be known by \emph{everyone}, but $a$ does not
        know $b$), so $b$ is eliminated.
\end{itemize}

Either way, a single call eliminates exactly one of the two candidates --
never both, and never neither -- so repeatedly applying this to pairs of
remaining candidates narrows the field down to a single survivor using
only $n-1$ calls total.

\begin{ideabox}[Stack-Driven Elimination, Then One Verification Pass]
Push every person $0, \dots, n-1$ onto a stack. While more than one
candidate remains: pop two, $a$ and $b$. If $\textit{knows}(a,b)$: push
$b$ back (eliminate $a$). Else: push $a$ back (eliminate $b$). After this
process, exactly one candidate remains -- but it is only a
\emph{candidate}, not yet confirmed (it is possible that \emph{no}
celebrity exists at all). Verify the survivor by checking that
\emph{everyone else} knows them and that they know \emph{no one else} --
an additional $O(n)$ scan.
\end{ideabox}

\subsection*{Why elimination narrows correctly, but still needs a final check}

The elimination process guarantees that \emph{if} a celebrity exists, they
are never eliminated: a true celebrity, compared against any other
candidate $x$, always satisfies $\textit{knows}(\textit{celebrity}, x) =
\text{false}$ (the celebrity knows no one) and
$\textit{knows}(x, \textit{celebrity}) = \text{true}$ (everyone knows the
celebrity) -- so whichever of the two checks the algorithm happens to make,
the celebrity is never the one eliminated. However, if \emph{no}
celebrity exists, the elimination process still produces \emph{some}
survivor (it always narrows to exactly one candidate, regardless of
whether that candidate has the celebrity property) -- which is why an
explicit final verification pass is required to confirm or reject the
survivor.

\subsection*{Algorithm}

\begin{enumerate}
  \item Push all $n$ people onto a stack.
  \item While the stack has more than one element:
  \begin{enumerate}
    \item Pop $a$, then pop $b$.
    \item If $\textit{knows}(a,b)$: push $b$. Else: push $a$.
  \end{enumerate}
  \item Let $\textit{candidate}$ be the single remaining person.
  \item Verify: for every other person $x$, check
        $\textit{knows}(x, \textit{candidate}) = \text{true}$ and
        $\textit{knows}(\textit{candidate}, x) = \text{false}$. If all
        checks pass, return $\textit{candidate}$; otherwise, return
        ``no celebrity exists.''
\end{enumerate}

\subsection*{Dry run}

\dryrun{Take $n=4$ with $\textit{knows}$ such that person $2$ is the true
celebrity: everyone knows $2$, and $2$ knows no one.}

Stack (top to bottom, after pushing $0,1,2,3$): top is $3$.

\begin{itemize}
  \item Pop $a=3$, $b=2$. Check $\textit{knows}(3,2)$: true (everyone
        knows $2$). Push $b=2$ (eliminate $3$). Stack $=[0,1,2]$.
  \item Pop $a=2$, $b=1$. Check $\textit{knows}(2,1)$: false ($2$ knows
        no one). Push $a=2$ (eliminate $1$). Stack $=[0,2]$.
  \item Pop $a=2$, $b=0$. Check $\textit{knows}(2,0)$: false. Push
        $a=2$ (eliminate $0$). Stack $=[2]$.
\end{itemize}

Candidate $=2$. Verification: every other person knows $2$ (true), and
$2$ knows none of them (true). Confirmed: person $2$ is the celebrity.

\subsection*{C++ implementation}

\begin{lstlisting}
#include <bits/stdc++.h>
using namespace std;

int findCelebrity(int n, function<bool(int,int)> knows) {
    stack<int> st;
    for (int i = 0; i < n; i++) st.push(i);

    while (st.size() > 1) {
        int a = st.top(); st.pop();
        int b = st.top(); st.pop();

        if (knows(a, b)) st.push(b); // eliminate a
        else st.push(a);             // eliminate b
    }

    int candidate = st.top();
    for (int x = 0; x < n; x++) {
        if (x == candidate) continue;
        if (!knows(x, candidate) || knows(candidate, x)) {
            return -1; // no celebrity
        }
    }
    return candidate;
}

int main() {
    // Person 2 is the celebrity: everyone knows 2, 2 knows no one.
    vector<vector<bool>> M = {
        {0,1,1,1},
        {0,0,1,0},
        {0,0,0,0},
        {0,1,1,0}
    };
    auto knows = [&](int a, int b) { return M[a][b]; };

    cout << "Celebrity: " << findCelebrity(4, knows) << endl;
    return 0;
}
\end{lstlisting}

\begin{complexitybox}
\textbf{Time Complexity.} The elimination phase makes exactly $n-1$
calls to $\textit{knows}$ (one per pop-pair, each removing one candidate
from a starting pool of $n$), and the verification phase makes up to
$2(n-1)$ more calls, giving an overall time complexity of
\[
  O(n).
\]
\textbf{Space Complexity.} $O(n)$ for the stack.
\end{complexitybox}

\begin{takeawaybox}
When a problem guarantees that comparing any two candidates always lets
you safely eliminate \emph{at least one} of them (never both, never
neither), a stack-driven elimination tournament narrows $n$ candidates
to $1$ in $O(n)$ comparisons. This is a different flavor of stack usage
from the monotonic-stack family -- here the stack just manages an
elimination order, not a sorted skyline -- and it is always worth pairing
with an explicit final verification step whenever the narrowing process
does not, by itself, guarantee the survivor actually has the desired
property.
\end{takeawaybox}
